hubi{This paper is not particularly well written if you ask me. Also, it doesn't have any citation (according to Google Scholar) from TCS, only in other fields like signal processing, biotechnology etc. The only reasonable reference is \cite{schleif_practical_2024}, though I cannot access it.}

This is the paper that inspired Monaldo's idea to have efficient ways to check if a matrix is PSD, by using a shift: if you have a matrix $A$ s.t. you want to know if it is PSD, you just compute its eigenvalues. 

Computing the eigenvalues, is the same as computing the values $\lambda$ s.t. $\det(A - \lambda I) = 0$. This generally takes $O(n^3)$.

A faster way to estimate the lower bound of the eigenvalues, is by using Gershgorin circles: define 

\begin{equation}
	R_i(A) := \sum_{j\neq i} |a_{i,j}|
\end{equation}
\begin{equation}
	D_i(A) := \{z \in \mathbb{C} : |a_{i,i} - z| < R_i(A)\}
\end{equation}

If $A$ is symmetric, $D_i(A) \subseteq \mathbb{R}$. In easier terms, we are just measuring the diagonal dominance of every row $i$.

\begin{theorem}[Gershgorin Circles]
\label{thm:gershgorin_circle}
	Spectrum$(A) \subseteq \cup_i D_i(A)$.
\end{theorem}

So any lower bound on the the Gershgorin circles is a good lower bound for the spectrum of $A$.



